% IW49091C
% Nadelstichverletzung
% 2.0.1
% 2015-11-15

.. _m-nadelstichverletzung:


\Z{Basiernd auf:cite:`OeAG2007`\ .}
\begin{itemize}
    -   Erstmaßnahmen durchführen: \nocite{OeAG2007}
    \begin{itemize}
        -   bei Stich- bzw.Schnittverletzungen:
        Blutung durch Kompression des umgebenden
        Gewebes induzieren. Danach desinfizieren und einen
        in Desinfektionsmittel getränkten Tupfer 10\Min*
        auf der Wunde belassen.
        -   bei Spritzern mit Blut oder Sekreten in Auge oder Mund:
        Schleimhaut mit Wasser u./o. Schleimhautdesinfektionsmittel spülen
    \end{itemize}
    -   Leitstelle informieren
    -   → **Umgehend in ein
    geeignetes Krankenhaus** (Von der Leitstelle erfragen).
    -   Patient: Zielspital informieren:
    \begin{itemize}
        -   Blut sichern für \Abk{HIV}- und Hepatitis-Serologie
    \end{itemize}
    -    **Unfallmeldung**  (wichtig!), Blutabnahme
    -   ggf. \Abk{HIV}-Postexpositionsprophylaxe (\Abk{PEP})
    \begin{itemize}
        -   kurzfristige Therapie (4 Wochen)
        -   Einzelfall muss immer individuell beurteilt werden
        (Nutzen/Risiko), starke Nebenwirkungen.
        -   Muss innerhalb von 2\Hour* begonnen werden → daher
        \underline{**sofort**} in ein geeignetes Krankenhaus!
    \end{itemize}
    -   Meldung an Betriebsleitung
    -   \EE{Unfallmeldung (wichtig!) an die \Abk{AUVA}}
\end{itemize}

\begin{Danger}
  -   Unfallmeldung nicht vergessen!
  
  Sofort in ein geeignetes Krankenhaus!
\end{Danger}
